\section{Conclusão}

Os experimentos permitiram observar o comportamento de um oscilador massa-mola vertical em oscilação livre, amortecida e forçada. No ar, o período foi $T_0=(1{,}399\pm0{,}005)\un{s}$, resultando em $\omega_0=(4{,}493\pm0{,}016)\un{rad/s}$ e, com a massa de $(10{,}84\pm0{,}12)\un{g}$, em uma constante elástica $k=(0{,}219\pm0{,}002)\un{N/m}$. Na água, o período aumentou para $T_1=(1{,}432\pm0{,}018)\un{s}$ e a frequência diminuiu para $\omega_1=(4{,}388\pm0{,}054)\un{rad/s}$, diferença no sentido previsto pelo modelo do oscilador amortecido.

As amplitudes na água apresentaram decaimento aproximadamente exponencial, com $\gamma=(0{,}37\pm0{,}03)\un{rad/s}$, evidenciando também a dissipação da energia mecânica pelo atrito. Esse valor, entretanto, não é consistente com o obtido pela relação $\omega_1=\sqrt{\omega_0^2-\gamma^2}$, de $0{,}96\un{rad/s}$, discrepância atribuída principalmente ao aumento da massa efetiva do corpo na água, que reduz a frequência natural em relação à medida no ar, à sensibilidade dessa relação a pequenos erros de cronometragem e ao caráter não puramente linear do atrito.

Na oscilação forçada, a frequência de ressonância no ar, $\Omega_R=(4{,}520\pm0{,}007)\un{rad/s}$, foi muito próxima da frequência natural $\omega_0$ (diferença de $0{,}6\%$), com amplitude muito grande. Na água, a ressonância ocorreu em $\Omega_R=(4{,}21\pm0{,}04)\un{rad/s}$, com amplitude máxima de $22{,}9\un{cm}$, menor e com variação mais suave que no ar, confirmando qualitativamente o efeito do amortecimento sobre a forma da curva e sobre a posição do máximo, previsto por $\Omega_r=\sqrt{\omega_0^2-2\gamma^2}$. Quantitativamente, o deslocamento medido foi maior que o previsto com o $\gamma$ obtido do decaimento, o que reforça que o modelo de atrito viscoso linear é apenas aproximado para o sistema na água.

Assim, a prática permitiu compreender experimentalmente os fenômenos de amortecimento e ressonância e evidenciou a importância da precisão das medidas, sobretudo as de cronometragem e de leitura das amplitudes, para a confiabilidade dos resultados.
